% LaTeX Curriculum Vitae Template
%
% Copyright (C) 2004-2009 Jason Blevins <jrblevin@sdf.lonestar.org>
% http://jblevins.org/projects/cv-template/
%
% You may use use this document as a template to create your own CV
% and you may redistribute the source code freely. No attribution is
% required in any resulting documents. I do ask that you please leave
% this notice and the above URL in the source code if you choose to
% redistribute this file.

\documentclass[letterpaper]{article}
\usepackage{CJKutf8}
\usepackage{hyperref}
\usepackage{geometry}
\usepackage{enumitem}
\usepackage{array}
% Comment the following lines to use the default Computer Modern font
% instead of the Palatino font provided by the mathpazo package.
%\usepackage[T1]{fontenc}
%\usepackage[sc,osf]{mathpazo}

% \hypersetup is kept outside any CJK block so keyval options parse correctly.
% PDF properties use Romanized metadata so they build without an active CJK block.
\hypersetup{
  colorlinks = true,
  urlcolor  = blue,
  pdfborder = {0 0 0},
  pdfauthor = {Zhiyuan Chen},
  pdfkeywords = {industrial organization, innovation, development, international trade},
  pdftitle = {Zhiyuan Chen - Curriculum Vitae},
  pdfsubject = {Curriculum Vitae},
  pdfpagemode = UseNone
}

\geometry{
  body={6.5in, 8.5in},
  left=1.25in,
  right=1.25in,
  top=1.2in
}

% Customize page headers
\pagestyle{myheadings}
\thispagestyle{empty}

% Custom section fonts
\usepackage{sectsty}
\sectionfont{\rmfamily\mdseries\Large}
\subsectionfont{\rmfamily\mdseries\itshape\large}

% Don't indent paragraphs.
\setlength\parindent{0em}

% Numbered paper list: [1], [2], ...
\newlist{papers}{enumerate}{1}
\setlist[papers]{label={[\arabic*]}, leftmargin=2.6em, itemsep=0.55em, topsep=0.35em}

\begin{document}
\begin{CJK}{UTF8}{gbsn}

% Redefine inside the document CJK block so tokens stay valid
\def\name{陈志远}
\markright{\name}

% Print name centered and bold:
\centerline{\huge \bf \name}

\vspace{1ex}

\begin{center}
{\scriptsize 最近更新：2026年9月20日}
\end{center}

\vspace{0.2in}

\normalsize

\begin{minipage}{0.55\linewidth}
  助理教授 \\
  \textbf{中国人民大学}商学院 \\
  北京市海淀区中关村大街59号 \\
  北京 100872
\end{minipage}
\begin{minipage}{0.45\linewidth}
  \begin{tabular}{ll}
    国籍： & 中国 \\
    电子邮箱： & \href{mailto:chenzhiyuan@rmbs.ruc.edu.cn}{\tt chenzhiyuan@rmbs.ruc.edu.cn} \\
    个人网页： & \href{https://zhiyuanryanchen.github.io/}{\tt zhiyuanryanchen.github.io} \\
  \end{tabular}
\end{minipage}

\section*{\textsc{工作经历}}
\textbf{中国人民大学}商学院 \\
企业经济学助理教授，2020年至今

\section*{\textsc{教育背景}}
\begin{minipage}{0.7\textwidth}
 经济学博士，\textbf{宾州州立大学}\\
 经济学硕士，\textbf{中国人民大学}  \\
 经济学学士（优等），\textbf{中国人民大学}
\end{minipage}
\begin{minipage}{0.3\textwidth}
  \raggedleft
  2015--2020 \\
  2013--2017 \\
  2009--2013
\end{minipage}

\section*{\textsc{学术研究}}
\subsection*{\textit{研究方向}}
 产业组织；发展经济学；国际贸易 \\
 特别关注：技术进步、创新与产业政策。

\subsection*{\textit{英文期刊发表}}
\begin{papers}
\item ``Identifying Treatment Effects on Productivity: Theory with an Application to Firms' Production Digitalization,'' with Moyu Liao and Karl Schurter, \textbf{\textit{RAND Journal of Economics}}, forthcoming（即将刊出）.

\item ``CO$_{2}$ Emission Regulation and Generation Allocation with Heterogeneous Coal-fired Generators,'' with Rong Luo and Li Su, \textbf{\textit{Journal of Development Economics}}, 2026, Vol.\ 180, 103771.

\item ``A Cost--Benefit Analysis of R\&D and Patents: Firm-Level Evidence from China,'' with Jie Zhang and Yuan Zi, \textbf{\textit{European Economic Review}}, Vol.\ 133, April 2021, 103633.

\item ``Import and Innovation: Evidence from Chinese Firms,'' with Jie Zhang and Wenping Zheng, \textbf{\textit{European Economic Review}}, 2017 (94), 205--220.

\item ``Micro-to-Macro Gains from Regional Trade Agreements: Evidence from RCEP,'' with Zhihao Wang and Fuxin Zhai, \textit{The World Economy}, forthcoming（即将刊出）.

\item ``The Impact of Digitalization on Productivity: Evidence from Chinese Manufacturing Firms,'' with Hongxi Chen, Jingtao Yi, and Jinqiu He, \textit{Management and Organization Review}, forthcoming（即将刊出）.

\item ``Bank Competition, Defaulting Risks, and Innovation,'' with Mengbo Zhang, Aomeng Zhang, and Xinyang Li, \textit{Canadian Journal of Economics}, forthcoming（即将刊出）.

\item ``Trade Policy Uncertainty and Market Diversification by Risk-Averse Firms,'' with Zhihao Wang, Ting Zhu, Kejian Gu, and Ying Tang, \textit{China Economic Review}, 2025, Vol.\ 91, 102400.

\item ``How Does Banking Expansion Influence City Development and Synergy? A New Perspective from Government Debts,'' with Yuyuan Wen and Hao Yu, \textit{Cities}, 2025, 162, 105923.

\item ``Made and Created in China: The Role of Processing Trade,'' with Aksel Erbahar and Yuan Zi, \textit{Scandinavian Journal of Economics}, 2025, 127, 390--426.

\item ``Detecting Learning from Importing: The Effect of Importing Behavior on R\&D Expenditures of Chinese Firms,'' with Jingtao Yi and Jinchao Huang, \textit{The World Economy}, 2024, 47(7), 3244--3278.

\item ``Is Cracking Down on Corruption Really Good for the Economy? Firm-Level Evidence from China,'' with Xin Jin and Xu Xu, \textit{Journal of Law, Economics, and Organization}, 2021, 37(2), 314--357.

\item ``Types of Patents and Driving Forces behind the Patent Growth in China,'' with Jie Zhang, \textit{Economic Modelling}, 2019 (80), 294--302.

\item ``The Bank--Firm Relationship: Helping or Grabbing?'' with Yong Li and Jie Zhang, \textit{International Review of Economics \& Finance}, 2016 (42), 385--403.
\end{papers}

\subsection*{\textit{工作论文}}
\begin{papers}
\item ``Externalities and Vertical Bargaining in Two-Sided Retail Platforms,'' with Rong Luo and Mark Roberts.

\item ``International Cooperation of Trade and Industrial Policy,'' with Zhihao Wang and Wei Xiang.

\item ``The Uneven Reach of Skill-Biased Technology,'' with Xiangyu Feng and Zexuan Liu.

\item ``Do Local Industrial Policies Drive Battery Energy Storage Innovation?'' with Xiangyang Zhao, Fuxin Zhai, Lun Li, and Shengyu Tao.

\item ``Detecting the Causal Impact of Strategic Decisions on Firm Performance,'' with Han Jiang, Moyu Liao, and Hao Sun.

\item ``Financing Innovation with Innovation,'' with Minjie Deng, Min Fang, and Yunzhou Shang, forthcoming（即将刊出）.

\item ``Path-Length-Independent No-Regret Learning for Distributed Online Generalized Nash Games over Unbalanced Networks,'' with Zhenhua Deng and Yulan Qian.

\item ``User Rating and App Updates,'' with Jingtao Yi, Sihan Meng, Sali Li, and Noman Shaheer.

\item ``Learning and Scaling,'' with Zexuan Liu and Liang Chen, \textbf{\textit{Strategic Management Journal}} 返修.

\item ``Intellectual Property Financing, Innovation, and Aggregate TFP,'' with Guanliang Hu, \textbf{\textit{Journal of Financial Economics}} 重投.

\item ``Data-Driven Vertical Integration,'' with Zexuan Liu, Fuxin Zhai, and Yao-yu Chih, \textbf{\textit{Management Science}} 返修.

\item ``Identification of Multiproduct Production Functions: Non-Parametric Limits and a CES Solution,'' with Moyu Liao and Jie Zhang, \textit{Econometric Reviews} 重投.

\item ``How Heterogeneous Are Productivity Gains from Exporting?'' with Moyu Liao, Hao Yu, and Jinqiu He, \textit{Journal of Productivity Analysis} 返修.
\end{papers}

\subsection*{\textit{中文期刊论文（部分）}}
\begin{papers}
\item “中国专利质量的测度，变化机制与经济增长效应”，\textbf{《经济研究》}，2026年第3期。

\item “创新链与产业链融合下的产业政策”，\textbf{《经济研究》}，2024年第9期，154--172页。［封面文章］

\item “对外技术引进与中国本土企业自主创新”，\textbf{《经济研究》}，2020年第7期，92--105页。

\item “中国企业创新补贴政策绩效评估：理论与证据”，\textbf{《经济研究》}，2015年第10期，4--17页。［首篇文章］

\item “中国出口国内附加值的测算与变化机制”，\textbf{《经济研究》}，2013年第10期，124--137页。［封面文章］

\item “中国知识溢出本地化的行政边界效应研究”，\textbf{《中国工业经济》}，2025年第1期，81--99页。

\item “减税激励的有偏技术进步效应研究——来自固定资产加速折旧政策的证据”，\textbf{《统计研究》}，2026年第1期。

\item “高新技术企业认定政策与企业创新”，\textbf{《中国经济学》}，2024年第2期，92--118页。

\item “专利对企业报酬生产率的作用效应估计与分解”，\textbf{《计量经济学报》}，2023年，第3卷第1期，166--194页。

\item “制造业集聚、污染与绿色发展实践路径：基于空间溢出模型的研究”，\textbf{《统计研究》}，2022年第9期，46--61页。

\item “贸易政策不确定性与企业出口产品范围”，\textbf{《国际贸易问题》}，2022年第6期，90--105页。

\item “出口与生产率关系的新检验：中国经验”，\textbf{《世界经济》}，2016年第6期，54--76页。

\item “进口是否引致了出口：中国出口奇迹的微观解读”，\textbf{《世界经济》}，2014年第6期。［首篇文章］

\item “融资约束、融资渠道与企业研发投入”，\textbf{《世界经济》}，2012年第10期，66--90页。
\end{papers}

\subsection*{\textit{专著}}
\textbf{《生产函数的结构估计与应用》}，经济管理出版社，2026年6月，ISBN 978-7-5243-1075-4。

\section*{\textsc{科研项目}}

\begin{tabular}{@{}p{0.80\textwidth}>{\raggedleft\arraybackslash}p{0.20\textwidth}@{}}
  国家自然科学基金面上项目：“数字经济下生产函数识别与技能偏向型技术进步测度”，主持人。 & 2027--2030 \\[0.9em]
  国家自然科学基金专项项目：“基础研究支撑经济发展的数据体系”，子课题负责人。 & 2026-- \\[0.9em]
  国家自然科学基金青年科学基金项目：“基于专利文本大数据的中国专利质量测度及其演化机制研究”，主持人。 & 2022--2024 \\[0.9em]
  腾讯研究院课题：“平台专利价值、创新与技术领导力”，主持人。 & 2023--2026 \\[0.9em]
  中国人民大学科学研究基金：“专利融资的量化评估与政策实践”，主持人。 & 2023--2025 \\[0.9em]
  中国人民大学哲学社会科学数据项目：“中国专利质量与价值数据库”，主持人。 & 2023 \\[0.9em]
  国家社会科学基金重大项目：“加快现代化经济体系建设”，子课题负责人。 & 2022 \\[0.9em]
  国家自然科学基金重点项目：“政府与市场功能协同的中国创新激励机制”，子课题负责人。 & 2021
\end{tabular}

\section*{\textsc{学术活动}}
\subsection*{中国人民大学教学}
\begin{minipage}{0.64\textwidth}
  经济与商业实证方法（硕博） \\
  机器学习与因果推断（本科） \\
  创新经济学（本科） \\
  产业政策专题（研究生） \\
  管理经济学，iMBA（英文授课） \\
  管理经济学，MBA（中文授课） \\
  全球网络周（GNAM）项目 \\
  市场分析与研究（本科）
\end{minipage}\hspace{0.02\textwidth}
\begin{minipage}{0.34\textwidth}
  \raggedleft
  2020年秋；2022--2026年春 \\
  2026年春 \\
  2023、2024年秋；2025、2026年春 \\
  2022--2024年春 \\
  2022--2024年秋；2025年春 \\
  2023、2024年秋；2025年春 \\
  2024--2026年春 \\
  2021年春
\end{minipage}

\subsection*{宾州州立大学教学}
\begin{minipage}{0.8\textwidth}
  \textbf{讲师}，宏观经济学原理（Econ 104） \\
  \textbf{讲师}，计量经济学导论（Econ 306） \\
  \textbf{助教}，中级微观经济学 \\
  \textbf{暑期研究实习生}，墨西哥中央银行
\end{minipage}
\begin{minipage}{0.2\textwidth}
  \raggedleft
  2019年夏 \\
  2018年夏 \\
  2017--2019 \\
  2017/07--2017/09
\end{minipage}

\subsection*{其他教学材料}
教学案例，发表于中国人民大学商学院案例中心，2023年、2025年。 \\
在编教材：\textit{AI辅助因果推断与机器学习}；\textit{创新经济学：理论与政策}，与张杰合著。

\subsection*{期刊审稿}
{\it Abacus, Applied Economics Letters, Asian Journal of Technology Innovation, China Economic Review, China \& World Economy, Economic Modelling, Economics Letters, Emerging Markets Finance and Trade, International Economics and Economic Policy, International Journal of Industrial Organization, International Review of Economics and Finance, Journal of Comparative Economics, Journal of Development Economics, Journal of International Management, Journal of Management Science and Engineering, Review of World Economy, The World Economy.}

中文期刊审稿：《经济研究》《中国社会科学》《金融研究》等；国际商务学会（AIB）年会会议审稿人，2024年。

\subsection*{会议、工作坊与邀请讲座}
{\footnotesize * 合作者报告。内部读书会从略。}

\vspace{0.6em}
\textbf{2026年：}
\begin{itemize}[leftmargin=1.6em, itemsep=0.25em, topsep=0.2em]
  \item “AI赋能一站式科研”邀请讲座系列：中国人民大学经济学院（6/4）；中央财经大学国际经济与贸易学院（6/9，7/7--8）；南京大学经济学院（6/15）；对外经济贸易大学WTO研究院（6/17）；首都经济贸易大学金融学院（6/24）与国际经济管理学院（6/24）；北京师范大学经济与工商管理学院（7/9）；深圳大学中国经济特区研究中心（7/14）。
  \item YES暑期学术讲座，西北大学，7/17。
  \item 第十届CCER暑期学院，北京大学国家发展研究院（崇礼），7/3--5。
  \item 中国数字经济发展与治理论坛年会，清华大学，7/4。
  \item 中国金融研究论坛（CFRC），清华大学五道口金融学院，7/2。
  \item 中国经济研究讲座与研讨会（第一期），特邀演讲，中国人民大学经济学院，1/20。
\end{itemize}

\textbf{2025年：}
\begin{itemize}[leftmargin=1.6em, itemsep=0.25em, topsep=0.2em]
  \item 第四届开放经济与宏观金融论坛，上海财经大学金融学院，12/27。
  \item 第二届香樟学堂学术讲座，主旨演讲，湖南科技大学商学院，12/19。
  \item 邀请报告，湖南师范大学商学院，12/18。
  \item 第五届香樟产业经济学论坛，主旨演讲，安徽工业大学，11/15。
  \item 海峡两岸微观经济学工作坊：理论与应用，厦门大学，9/27--28。
  \item 邀请报告，澳门大学社会科学学院，9/3。
  \item Eaton论坛，国际经济与金融学会（IEFS）中国分会，对外经济贸易大学，8/14。
  \item 青年学者学术沙龙暨第二届青年发展经济学论坛，山东大学，7/7--8。
  \item 第三届全球校友学术年会，中国人民大学财政金融学院，7/4。
  \item ISEM会议：中国实证研究前沿，首都经济贸易大学，5/19。
\end{itemize}

\textbf{2024年：}
\begin{itemize}[leftmargin=1.6em, itemsep=0.25em, topsep=0.2em]
  \item 工商管理学科建设中国自主知识体系联盟论坛，“新质生产力与数字经济”分论坛；发布《平台企业数字技术发明专利研究报告》，中国人民大学商学院，4/21。
\end{itemize}

\textbf{2023年：}
\begin{itemize}[leftmargin=1.6em, itemsep=0.25em, topsep=0.2em]
  \item 宾州州立大学校友实证产业组织工作坊，中国人民大学经济学院，7/6。
  \item 中国金融研究论坛，清华大学五道口金融学院，7/2。
  \item 邀请报告，南京大学商学院，4/7。
\end{itemize}

\textbf{2022年：}
\begin{itemize}[leftmargin=1.6em, itemsep=0.25em, topsep=0.2em]
  \item 邀请报告，中国人民大学国家发展与战略研究院，12/14。
  \item 邀请报告，对外经济贸易大学全球价值链研究院，10/27。
  \item 邀请报告，复旦大学中国社会主义市场经济研究中心，10/11。
  \item 第九届新结构经济学国际研讨会，评论人，北京大学新结构经济学研究院，8/2。
  \item 第六届CCER暑期学院，北京大学国家发展研究院，7/2。
  \item 2022年世界计量经济学会亚洲年会（AMES），香港中文大学（深圳），6/22。
\end{itemize}

\textbf{2021年：}
\begin{itemize}[leftmargin=1.6em, itemsep=0.25em, topsep=0.2em]
  \item 邀请报告，北京大学经济学院，11/19。
  \item 第七届国际青年金融学者会议，北京大学汇丰商学院，7/12。
\end{itemize}

\textbf{2020年：}
\begin{itemize}[leftmargin=1.6em, itemsep=0.25em, topsep=0.2em]
  \item 邀请报告，中国人民大学商学院，9/28。
\end{itemize}

\textbf{2019年：}
\begin{itemize}[leftmargin=1.6em, itemsep=0.25em, topsep=0.2em]
  \item 邀请报告：对外经济贸易大学国际经济贸易学院（12/12）、湘潭大学商学院（12/10）、湖南大学经济与贸易学院（12/4）、香港中文大学（深圳）、北京大学新结构经济学研究院、中国人民大学经济学院。
  \item 第14届经济学研究生年会，圣路易斯华盛顿大学。
  \item *中西部国际经济发展会议。
\end{itemize}

\textbf{2018年：}
\begin{itemize}[leftmargin=1.6em, itemsep=0.25em, topsep=0.2em]
  \item *首届CICE年会，清华大学。
  \item 美国经济学会年会（海报展示），费城。
  \item *应用与理论经济学工作坊（WATE）。
  \item *“中国创新与增长的管理与建模挑战”特刊会议，中国人民大学。
\end{itemize}

\textbf{2017年：}
\begin{itemize}[leftmargin=1.6em, itemsep=0.25em, topsep=0.2em]
  \item 世界计量经济学会中国年会，武汉。
  \item *中国研究国际联盟年会，曼海姆。
\end{itemize}

\textbf{2014年：} 日内瓦——中国年度工作坊，日内瓦。

\vspace{0.3em}
\textbf{2013年：} 复旦大学博士生论坛。

\section*{\textsc{荣誉奖励}}

\textit{国家青年人才计划}，2025年

\vspace{0.5em}
\textit{中国知网经济学高被引学者（前1\%）}，2025年、2024年

\vspace{0.5em}
\textit{第七届张培刚发展经济学优秀论文奖}，张培刚发展经济学研究基金会，2018年

\vspace{0.5em}
\textit{第七届高等学校科学研究优秀成果奖（人文社会科学）三等奖}，教育部

\vspace{0.5em}
\textit{2016年国际贸易学十佳中文论文}，中国社会科学院世界经济与政治研究所，2017年

\vspace{0.5em}
\textit{第18届安子介国际贸易研究奖三等奖}，对外经济贸易大学，2014年

\vspace{0.5em}
\textit{国家奖学金（前1\%）}，中国人民大学，2012年、2013年、2014年

\section*{\textsc{专利、软件与技能}}
\textbf{专利：}“基于专利文本相似度的专利质量测算方法及系统”（实质审查中）。

\vspace{0.5em}
\textbf{软件：}Stata模块\texttt{pte}（生产率处理效应估计）作者之一（与蔡晓慧、廖漠雨、徐伟合作），2026年。熟练使用MATLAB、Stata、Julia、Python。

\vspace{0.5em}
\textbf{语言：}中文（母语），英文（流利）

\clearpage
\end{CJK}
\end{document}
